\documentclass[12pt,letterpaper]{article}

% ==========================================================
% PLANTILLA LATEX / ICONTEC
% Curso Metodología de la Investigación - UPTC
% Generada a partir del Instrumento de trabajo (Lección 1):
% "Selección de la idea de investigación" - únicamente la
% información propia del equipo (no el ejemplo resuelto).
% ==========================================================

\usepackage[top=3cm,left=4cm,bottom=3cm,right=2cm]{geometry}
\usepackage[utf8]{inputenc}
\usepackage[T1]{fontenc}
\usepackage{mathptmx} % tipografía tipo Times, similar al PDF de referencia
\usepackage{setspace}
\usepackage{titlesec}
\usepackage{tabularx}
\usepackage{longtable}
\usepackage{booktabs}
\usepackage{array}
\usepackage[table]{xcolor}
\usepackage[most]{tcolorbox}
\usepackage{enumitem}
\usepackage{hyperref}
\hypersetup{colorlinks=false, hidelinks}

\onehalfspacing

\definecolor{badgebg}{HTML}{FDF6E3}
\definecolor{badgeborder}{HTML}{C9A227}
\definecolor{badgetext}{HTML}{8A6D1D}
\definecolor{greensignal}{HTML}{2E7D32}

% ---------- Formato de títulos de sección (centrado, versalitas, como el PDF) ----------
\titleformat{\section}{\centering\bfseries\large}{\thesection.}{0.5em}{\MakeUppercase}
\titleformat{\subsection}{\bfseries\normalsize}{\thesubsection}{0.5em}{}
\titlespacing*{\section}{0pt}{2.5em}{1.5em}

% ---------- Comando para el "badge" EN CONSTRUCCIÓN / estado ----------
\newcommand{\estado}[1]{%
  \begin{center}
  \tcbox[colback=badgebg,colframe=badgeborder,boxrule=0.5pt,
    arc=6pt,left=6pt,right=6pt,top=3pt,bottom=3pt,
    fontupper=\small\bfseries\color{badgetext}]{\MakeUppercase{#1}}
  \end{center}
}

% ---------- Comando para la nota itálica que acompaña cada estado ----------
\newcommand{\notafase}[1]{\begin{center}\itshape\small #1\end{center}}

\newcolumntype{L}[1]{>{\raggedright\arraybackslash}p{#1}}

\begin{document}

% ================= PORTADA =================
\begin{titlepage}
\begin{center}
\vspace*{1.5cm}
{\large UNIVERSIDAD PEDAGÓGICA Y TECNOLÓGICA DE COLOMBIA}\\[0.3cm]
{\large FACULTAD DE INGENIERÍA}\\[0.3cm]
{\large INGENIERÍA DE SISTEMAS Y COMPUTACIÓN}

\vspace{3cm}

{\Large\bfseries RELACIÓN ENTRE LAS HORAS DE EXPOSICIÓN A PANTALLAS,\\
LOS HÁBITOS POSTURALES Y LOS SÍNTOMAS FÍSICOS EN\\
ESTUDIANTES DE INGENIERÍA DE SISTEMAS DE LA UPTC}

\vspace{0.8cm}
{\itshape\small Título provisional --- se ajusta al formular la propuesta (Fase 3)}

\vspace{3cm}

{\bfseries AUTORES}\\[0.6cm]
Gabriel Alejandro Gamba Leguizamón\\
José Fernando Barreto Franco\\
Daniel Felipe Molina Barajas

\vfill

Documento elaborado en el marco del curso Metodología de la Investigación\\[0.3cm]
Tunja, Colombia · 2026

\end{center}
\end{titlepage}

% ================= CONTENIDO (tabla de estado por sección) =================
\section*{\centering CONTENIDO}
\begin{center}
\begin{tabularx}{\textwidth}{L{10cm} X}
\toprule
\textbf{Sección} & \textbf{Estado}\\
\midrule
Índice de tablas & incluido\\
Índice de figuras & incluido\\
1. Planteamiento del problema & parcial\\
\quad 1.1 Antecedentes & en construcción\\
\quad 1.2 Descripción de la situación problemática & redactado\\
\quad 1.3 Formulación del problema & en construcción\\
\quad 1.4 Consecuencias de no investigar & en construcción\\
\quad 1.5 Delimitación & en construcción\\
2. Justificación & en construcción\\
3. Objetivos & en construcción\\
4. Alcance de la investigación & en construcción\\
5. Formulación de hipótesis & en construcción\\
6. Variables & en construcción\\
7. Marco referencial & en construcción\\
8. Diseño metodológico & en construcción\\
9. Resultados y discusión & en construcción\\
10. Conclusiones y recomendaciones & en construcción\\
Referencias & en construcción\\
Anexos & incluido (parcial)\\
\bottomrule
\end{tabularx}
\end{center}

\newpage
% ================= ÍNDICE DE TABLAS =================
\section*{\centering ÍNDICE DE TABLAS}
Tabla 1. Causas y síntomas identificados en la situación problemática\dotfill 1

\vspace{1cm}
\notafase{Esta lista crece con cada tabla que se incorpore al cuerpo del documento en fases posteriores.}

\newpage
% ================= ÍNDICE DE FIGURAS =================
\section*{\centering ÍNDICE DE FIGURAS}
\notafase{Aún no hay figuras incorporadas al cuerpo del documento. Se listan aquí a medida que cada fase incorpore figuras propias (diagramas, gráficos, esquemas, etc.).}

\newpage
% ================= 1. PLANTEAMIENTO DEL PROBLEMA =================
\section{Planteamiento del problema}

\subsection{Antecedentes}
\estado{En construcción}
\notafase{Se redacta cuando se recopile la evidencia y las cifras del contexto (Fase 2 del curso).}

\subsection{Descripción de la situación problemática}

La problemática principal de esta investigación surge del impacto que generan las
extensas jornadas de trabajo frente al computador sobre la salud física y el
bienestar general de los estudiantes de Ingeniería de Sistemas. El uso continuo e
ininterrumpido de pantallas, sumado a la adopción de posturas inadecuadas durante
las sesiones de programación y estudio, se asocia con molestias osteomusculares y
afecciones visuales que pueden condicionar el rendimiento académico, la
concentración y la calidad de vida de los futuros ingenieros, quienes a menudo
normalizan el malestar físico como un aspecto inevitable de su disciplina.

\medskip
\noindent\textbf{Causas.} Entre las principales causas se encuentra el
desconocimiento de las normas básicas de higiene postural y de los principios de
ergonomía en entornos de aprendizaje digital, así como el uso de mobiliario y
equipos no adaptados ergonómicamente (por ejemplo, computadores portátiles usados
sin soporte elevado ni periféricos externos, en mesas o sillas no regulables).
También influyen la falta de hábito en la realización de pausas activas --- favorecida
por la alta exigencia de entregas de software que fomenta jornadas continuas sin
descansos--- y condiciones ambientales desfavorables, como una iluminación deficiente,
distancias inadecuadas frente a la pantalla y un ajuste de brillo incorrecto.

\medskip
\noindent\textbf{Síntomas.} Entre los síntomas asociados se observan episodios
recurrentes de fatiga visual, ardor ocular, sequedad y visión borrosa (vinculados al
Síndrome Visual Informático), así como tensiones y dolores crónicos en la zona
cervical, lumbar, hombros y muñecas derivados de mantener posturas forzadas por
periodos prolongados. Se suma la fatiga física y mental acumulada: la permanencia
continuada frente a dispositivos altera los patrones de sueño, reduce los niveles de
atención y la capacidad de concentración, y el sedentarismo asociado contribuye al
descondicionamiento físico general y a la aparición de cefaleas al finalizar las
actividades diarias.

\bigskip
\begin{center}
\begin{longtable}{L{7cm} L{7cm}}
\caption{Causas y síntomas identificados en la situación problemática} \\
\toprule
\textbf{Causas} & \textbf{Síntomas} \\
\midrule
\endfirsthead
Desconocimiento de normas ergonómicas & Fatiga visual (Síndrome Visual Informático) \\
Uso de mobiliario y equipos inadecuados & Dolores osteomusculares (cervical, lumbar, muñecas) \\
Ausencia de pausas activas durante la jornada & Cansancio físico y mental acumulado \\
Exposición ininterrumpida a pantallas & Cefaleas y dolores de cabeza recurrentes \\
Iluminación y distancia de visión incorrectas & Disminución de la concentración y fatiga crónica \\
Estilo de vida sedentario y jornadas extensas & Alteraciones en los patrones de sueño \\
\bottomrule
\end{longtable}
\end{center}

\subsection{Formulación del problema}
\estado{En construcción}
\notafase{Se redacta cuando se libere el tema correspondiente en el curso (Fase 2).}

\subsection{Consecuencias de no investigar}
\estado{En construcción}
\notafase{Se redacta cuando se libere el tema correspondiente en el curso (Fase 2).}

\subsection{Delimitación}
\estado{En construcción}
\notafase{Se redacta cuando se libere el tema correspondiente en el curso (Fase 2). Delimitación tentativa: estudiantes de Ingeniería de Sistemas de la UPTC.}

\newpage
% ================= 2. JUSTIFICACIÓN =================
\section{Justificación}
\estado{En construcción}
\notafase{Se redacta en la Fase 3 del curso (Formulación de la propuesta). Insumo disponible en el Anexo A: respuestas de interés y pertinencia del instrumento de selección de idea.}

\newpage
% ================= 3. OBJETIVOS =================
\section{Objetivos}
\estado{En construcción}
\notafase{Se redacta en la Fase 3 del curso.}

\newpage
% ================= 4. ALCANCE =================
\section{Alcance de la investigación}
\estado{En construcción}
\notafase{Sección incorporada según Hernández-Sampieri (tipo de estudio: exploratorio, descriptivo, correlacional o explicativo) --- no aparece como título propio en la NTC 1486. Se redacta en la Fase 3.}

\newpage
% ================= 5. HIPÓTESIS =================
\section{Formulación de hipótesis}
\estado{En construcción}
\notafase{Sección incorporada según Hernández-Sampieri, Tamayo y Tamayo y Bernal --- no aparece como título propio en la NTC 1486. Se redacta en la Fase 3.}

\newpage
% ================= 6. VARIABLES =================
\section{Variables}
\estado{En construcción}
\notafase{Definición nominal/conceptual en la Fase 3; operacionalización completa en la Fase 5.}

\newpage
% ================= 7. MARCO REFERENCIAL =================
\section{Marco referencial}
\estado{En construcción}
\notafase{Se redacta en la Fase 4 del curso (Construcción del marco referencial).}

\newpage
% ================= 8. DISEÑO METODOLÓGICO =================
\section{Diseño metodológico}
\estado{En construcción}
\notafase{Se redacta en la Fase 5 del curso. Insumo disponible en el Anexo A: respuestas de viabilidad del instrumento de selección de idea (encuesta estructurada, participantes disponibles, herramientas de recolección y análisis).}

\newpage
% ================= 9. RESULTADOS =================
\section{Resultados y discusión}
\estado{En construcción}
\notafase{Se redacta en la Fase 6 del curso.}

\newpage
% ================= 10. CONCLUSIONES =================
\section{Conclusiones y recomendaciones}
\estado{En construcción}
\notafase{Se redacta en la Fase 7 del curso.}

\newpage
% ================= REFERENCIAS =================
\section*{\centering REFERENCIAS}
\estado{En construcción}
\notafase{Aún no se han incorporado referencias bibliográficas al proyecto. Se documentan aquí a medida que cada fase incorpore fuentes citadas, siguiendo la numeración de primera aparición en el texto.}

\newpage
% ================= ANEXOS =================
\section*{\centering ANEXOS}

\subsection*{Anexo A. Instrumento de selección de idea de investigación (Lección 1)}
\notafase{Reproduce, para trazabilidad, la información propia del equipo diligenciada en el instrumento de la Lección 1 (se excluye el ejemplo resuelto de la guía).}

\bigskip
\noindent\textbf{Datos del equipo}
\begin{center}
\begin{tabularx}{\textwidth}{L{5cm} X}
\toprule
\textbf{Campo} & \textbf{Respuesta}\\
\midrule
Integrante 1 & Gabriel Alejandro Gamba Leguizamón\\
Integrante 2 & José Fernando Barreto Franco\\
Integrante 3 & Daniel Felipe Molina Barajas\\
Fecha & 23 de agosto de 2026\\
Tema de interés general & Salud Digital, Ergonomía y Bienestar Físico en Estudiantes de Ingeniería de Sistemas\\
\bottomrule
\end{tabularx}
\end{center}

\bigskip
\noindent\textbf{Idea de investigación tentativa}

\noindent Quiero saber cuál es la relación entre el número de horas diarias
dedicadas a la exposición a pantallas, las posturas adoptadas durante las jornadas
de estudio y la frecuencia de síntomas físicos (fatiga visual y molestias
osteomusculares) en los estudiantes de Ingeniería de Sistemas de la UPTC, porque
eso importaría para proponer una guía práctica de ergonomía y pausas activas
adaptada a las rutinas de los futuros profesionales de software.

\bigskip
\noindent\textbf{Interés y pertinencia}
\begin{center}
\begin{tabularx}{\textwidth}{L{5.5cm} X}
\toprule
\textbf{Pregunta} & \textbf{Respuesta del equipo}\\
\midrule
¿Cuál es su interés personal en este tema? & Como estudiantes de Ingeniería de Sistemas, pasamos largas jornadas frente al computador y frecuentemente experimentamos fatiga visual o dolores de espalda y cuello, por lo que queremos entender cómo nuestras rutinas afectan la salud física.\\
¿Cuál es su interés académico o profesional? & El tema se relaciona directamente con la salud ocupacional e higiene postural en el ámbito de la ingeniería de software, permitiéndonos adquirir conocimientos sobre la prevención de riesgos físicos derivados del trabajo sedentario y continuo con computadores.\\
¿Es un tema prioritario o significativo para la disciplina? & Sí. La extensión de jornadas de programación y clases virtuales aumenta la incidencia de síndromes como el Síndrome Visual Informático y trastornos musculoesqueléticos en estudiantes e ingenieros, impactando su rendimiento y calidad de vida.\\
¿A quién afecta el problema que identificaron? & Principalmente a los estudiantes de la Escuela de Ingeniería de Sistemas, docentes, y profesionales del área que pasan más de 6 horas diarias expuestos a pantallas.\\
\bottomrule
\end{tabularx}
\end{center}

\bigskip
\noindent\textbf{Originalidad}
\begin{center}
\begin{tabularx}{\textwidth}{L{5.5cm} X}
\toprule
\textbf{Pregunta} & \textbf{Respuesta del equipo}\\
\midrule
¿Ya se ha investigado esto antes? & Existen estudios generales sobre ergonomía y uso de pantallas en trabajadores de oficina, pero pocos enfocados específicamente en la población universitaria de Ingeniería de Sistemas de nuestro contexto regional.\\
Si ya existe, ¿qué preguntas quedan sin responder? & Cómo se relacionan específicamente el número acumulado de horas continuas de programación con la aparición temprana de molestias osteomusculares y visuales en estudiantes jóvenes.\\
¿Qué aportarían ustedes que no exista ya? & Datos primarios y actualizados del contexto universitario local (UPTC) junto con un análisis correlacional sencillo entre hábitos posturales, horas de pantalla y síntomas físicos reportados.\\
\bottomrule
\end{tabularx}
\end{center}

\bigskip
\noindent\textbf{Viabilidad}
\begin{center}
\begin{tabularx}{\textwidth}{L{5.5cm} X}
\toprule
\textbf{Pregunta} & \textbf{Respuesta del equipo}\\
\midrule
¿Qué tan complejo es el tema para desarrollarlo en 16 semanas? & Muy viable y acotado. Consiste en diseñar y aplicar una encuesta estructurada (usando escalas validadas como la escala de usabilidad o dolor sintomático) y analizar los datos correlacionales.\\
¿Tienen acceso a los datos, herramientas o participantes que necesitarían? & Sí. Tenemos acceso directo a nuestros compañeros de carrera como participantes del estudio y utilizaremos herramientas gratuitas como Google Forms para la recolección de datos y programas como Excel para el análisis de frecuencia.\\
¿Tienen (o pueden adquirir en el curso) los conocimientos técnicos necesarios? & Sí. El diseño de encuestas y el análisis estadístico descriptivo/correlacional básico se abordan en el marco del curso de Metodología de la Investigación.\\
¿El tiempo disponible (4 h de clase + 4 h autónomas por semana) alcanza para esta idea? & Sí. El tiempo es suficiente para la revisión de antecedentes, diseño del cuestionario, recolección de respuestas en 2 semanas, análisis estadístico y redacción del informe final.\\
\bottomrule
\end{tabularx}
\end{center}

\bigskip
\noindent\textbf{Semáforo de la idea}
\begin{center}
\begin{tabularx}{\textwidth}{L{2.5cm} L{6cm} X}
\toprule
\textbf{Señal} & \textbf{Significa} & \textbf{¿Aplica a la idea?}\\
\midrule
\textcolor{greensignal}{\textbf{Verde}} & Cumple los tres criterios sin ajustes mayores. & \textbf{Sí} --- puede avanzar a la Lección 2 (antecedentes).\\
Amarillo & Cumple en general, pero un criterio necesita ajustarse. & No\\
Rojo & Falla claramente en al menos un criterio. & No\\
\bottomrule
\end{tabularx}
\end{center}

\end{document}